% !PRINTED 329
\weberpage{329}
%<<CONT
Mal nach Zeilen, das andere Mal nach Columnen multiplicirt, so
folgt nach den Formeln (4):

\weq{5}{C_1\,(y,b) = r^6\,C_1\,(x,a),}

wodurch die Invarianten-Eigenschaft ausgedrückt und das Gewicht
$= 6$ bestimmt ist.

Die dritte Covariante $C_2$, die wir betrachten, ist die Functionaldeterminante
 aus den drei Formen $f$, $H$, $C_1$, also

\weq{6}{C_2\,(x, a) = \begin{vmatrix} f'\,(x_1), & H'\,(x_1), & C_1'\,(x_1) \\ f'\,(x_2), & H'\,(x_2), & C_1'\,(x_2) \\ f'\,(x_3), & H'\,(x_3), & C_1'\,(x_3) \end{vmatrix},}

die, wie schon aus \S.\,59 des ersten Bandes hervorgeht, die
Invarianten-Eigenschaft hat:

\weq{7}{C_2\,(y,b) = r^9\,C_2\,(x,a),}

und die also das Gewicht 9 hat.

Bezeichnen wir mit $\lambda$ das Gewicht, mit $\nu$ den Grad in den
Variablen, mit $\mu$ den Grad in den Coëfficienten, so ergiebt sich
folgende Zusammenstellung:

\begin{center}
\begin{tabular}{lccc}
 & $\nu$, & $\mu$, & $\lambda$ \\
\hline
$f$ & $n$, & 1, & 0 \\
$H$ & $3n-6$, & 3, & 2 \\
$C_1$ & $8n-18$, & 8, & 6 \\
$C_2$ & $12n-27$, & 12, & 9.
\end{tabular}
\end{center}

\begin{center}
\S.\,88.

\gesperrt{Die Hesse'sche Curve.}
\end{center}

Die erste der eingeführten Covarianten $H$ hat eine einfache
geometrische Bedeutung, die wir jetzt kennen lernen wollen.
Diese Untersuchung wird dadurch wesentlich erleichtert, dass
wir, wegen der Invarianten-Eigenschaft der Form $H$, nichts an
Allgemeinheit verlieren, wenn wir dem Coordinatensysteme irgend
eine specielle Lage gegen die Curve $f$ geben.

Legen wir die Ecke $\xi_3$ des Coordinatendreiecks in einen
Punkt der Curve $f$, der nicht zu den singulären gehört, und
wählen die Tangente in diesem Punkte zur $x_1$-Axe, so erhält
nach \S.\,86, (3) die Form $f$ den Ausdruck:

\weq{1}{f = h\,x_3^{n-1}x_1 + x_3^{n-2}(a\,x_1^2+2\,b\,x_1x_2+c\,x_2^2)+x_3^{n-3}f_3+\cdots,}
%CONT>>
